\documentclass[10pt,a4paper]{amsart}
\usepackage[margin=19mm]{geometry}
\usepackage{amsmath,amssymb,mathtools}
\usepackage[scr=rsfs,cal=euler]{mathalfa}
\usepackage{xfrac}
\usepackage[unicode,hidelinks,bookmarks=false]{hyperref}
\newcommand{\ie}{i.e.\ }
\newcommand{\cf}{cf.\ }
\newcommand{\resp}{respectively\ }
\newcommand{\Subsection}{\subsection}
\providecommand{\nobreakdash}{\nobreak\hbox{-}}
\setlength{\emergencystretch}{2em}
\multlinegap0pt
\usepackage{bm}
\usepackage{bbm}
\usepackage{dsfont}
\usepackage{xspace}
\usepackage{cancel}
\usepackage{mathtools}
\usepackage{amsthm}
\makeatletter
\@ifundefined{theorem}{\newtheorem{theorem}{Theorem}}{}
\@ifundefined{corollary}{\newtheorem{corollary}[theorem]{Corollary}}{}
\@ifundefined{lemma}{\newtheorem{lemma}[theorem]{Lemma}}{}
\@ifundefined{proposition}{\newtheorem{proposition}[theorem]{Proposition}}{}
\makeatother
\title[Additive structures imply more distances in $\mathbb{F}\_q^d$]{Additive structures imply more distances in $\mathbb{F}\_q^d$}
\author{Daewoong Cheong, Gennian Ge, Doowon Koh, Thang Pham, Dung The Tran, Tao Zhang}
\date{Source: arXiv:2510.26364v3 (CC BY 4.0)}
\begin{document}
\maketitle

\noindent\emph{Source and licence.} Additive structures imply more distances in $\mathbb{F}_q^d$. Version arXiv:2510.26364v3 (CC BY 4.0), distributed under the Creative Commons Attribution 4.0 International licence, as recorded in the arXiv licence metadata for this exact version and shown on the abstract page https://arxiv.org/abs/2510.26364v3. Full rights notice: licenses/arxiv-2510.26364-CC-BY-4.0.txt.\par\smallskip

\noindent\textit{Editorial scope.} Counted unit: the Proposition whose source label is \texttt{propro1} in the article (source0.tex lines 608--617, source label \texttt{propro1}), delivered together with the article's own complete original proof of it (source0.tex lines 619--668, one \texttt{proof} environment of 50 lines). No mathematics is written, repaired or re-derived: statement and proof are the article's own text line for line. Required context retained uncounted: cor:euivalence-Salem-bound (lines 426--438); SizePS (lines 522--522); lemma2.6 (lines 599--604). Every definition, display and label of those lines is retained unchanged. Omitted from this package and not redistributed: 1--425, 439--521, 523--598, 605--607, 618--618, 669--1494. Nothing inside a retained excerpt is omitted or altered: every retained body is delivered exactly as the article's lines are written, so no omission receipt of any kind accompanies this package. Citations: the retained excerpts carry no citation command at all, so no bibliography entry of the article is transcribed and no reference is redistributed. Reference adaptation: none was needed and none was made; every reference inside the delivered unit resolves inside it, so no unresolved reference or citation is produced. Printed slips kept literal: no printed word, symbol, hypothesis or number of the counted statement or of its proof was altered. Layout and class: the article's own class is \texttt{article} with packages including a4wide, amsmath, amssymb, amsthm, xcolor, parskip, hyperref, setspace; the wrapper is the lane's pinned \texttt{amsart} editorial wrapper at 10pt on A4, which restates the theorem-like environments the retained text needs and adds the article's own macro definitions that the retained text needs (none), with extra wrapper packages (bm, bbm, dsfont, xspace, cancel, mathtools). The delivered numbering is the unit's own. Assets: no retained span contains a figure environment, a graphics-inclusion command, a PDF-page inclusion, a table or a TikZ picture; no image file is copied into this package, the package declares no asset dependency and no image file is redistributed with it. Licence: version arXiv:2510.26364v3 of this article is distributed under the Creative Commons Attribution 4.0 International licence, as recorded in the arXiv licence metadata for this exact version and shown on the abstract page https://arxiv.org/abs/2510.26364v3. A full rights notice is delivered at licenses/arxiv-2510.26364-CC-BY-4.0.txt. Characters: every retained excerpt is pure ASCII, so the delivered engine, XeLaTeX with its default Latin Modern text font, reports no missing character.
\begin{corollary}\label{cor:euivalence-Salem-bound}
	For any $s\in[0,1]$, the $(2k, s)$--Salem condition
	\begin{equation*}
		\|\widehat{E}\|_{2k}
		\;\ll\; q^{-d }\,|E|^{1-s}
	\end{equation*}
	is equivalent to the additive energy bound
	\begin{equation*}
		\Lambda_{2k}(E)
		\;\ll\;
		\; \frac{|E|^{2k}}{q^d}+ |E|^{2k(1-s)}.
	\end{equation*}
\end{corollary}

\begin{equation}\label{SizePS} |P|\sim |S_j|\sim |S_g|\sim q^{d-1},\end{equation} 

\begin{lemma}\label{lemma2.6}
Suppose $d\ge 3$ is odd. For $E \subseteq S_j$ with $j\ne 0$, we have
\[
\Lambda_4(E) \ll \frac{|E|^3}{q} + q^{\frac{d-1}{2}}|E|^2.
\]
\end{lemma}

\begin{proposition}\label{propro1}
Assume that $E \subseteq S_j$, $j\ne 0$, with $d \ge 3$ odd, and that one of the following conditions holds
\begin{enumerate} 
    \item[(i)] $q^{\frac{d-1}{4(1-2s)}} \ll |E| \ll q^{\frac{d+1}{2}}, \qquad~~~\frac{1}{4}\leq s \leq \frac{1}{4}+\frac{1}{2(d+1)}$; 
    \item[(ii)] $q^{\frac{d+1}{2}} \ll |E| \ll q^{\,d-1}, \qquad ~~~~~~~ \dfrac{1}{4} \le s \le \dfrac{1}{4} + \dfrac{1}{4(d-1)}$;
    \item[(iii)] $q^{\frac{d+1}{2}} \ll |E| \ll q^{\frac{1}{4s-1}}$,\qquad~~~~~ $\frac{1}{4}+\frac{1}{4(d-1)} \leq s \leq \frac{1}{4}+\frac{1}{2(d+1)}$; 
    \item[(iv)] $|E|\sim q^{d-1} $,\qquad~~~~~~~~~~~~~~~~ $ \frac{1}{4} \leq s\le\frac{1}{2}$. 
\end{enumerate}
Then $E$ is a $(4,s)$\textnormal{-Salem set.}
\end{proposition}

\begin{proof}

By Corollary \ref{cor:euivalence-Salem-bound}, $E$ satisfies the $(4, s)$--Salem property precisely when
\[ \Lambda_4(E) \ll \frac{|E|^4}{q^d}+ |E|^{4-4s} .\]
Therefore, using Lemma  \eqref{lemma2.6}, our task reduces to verifying that under each of the assumptions $(i), (ii), (iii), (iv),$ the following holds
\begin{equation} \label{LR*} L:= \frac{|E|^3}{q} + q^{\frac{d-1}{2}}|E|^2 \ll \frac{|E|^4}{q^d}+ |E|^{4-4s}=: R. \end{equation}

\medskip
\noindent
\textbf{Case 1. Suppose $\frac{|E|^3}{q} \ll q^{\frac{d-1}{2}}|E|^2 (\iff |E| \ll q^{\frac{d+1}{2}})$}.
Under this assumption,  we see that $L\sim q^{\frac{d-1}{2}}|E|^2$. Hence,  \eqref{LR*} becomes
\begin{equation}\label{Suf*} q^{\frac{d-1}{2}}|E|^2 \ll \frac{|E|^4}{q^d}+ |E|^{4-4s}.\end{equation}
We establish that assumption $(i)$ guarantees this inequality.
Notice that  a sufficient condition for  \eqref{Suf*} is given by 
\[ q^{\frac{d-1}{2}}|E|^2 \ll \frac{|E|^4}{q^d}  \quad \mbox{or} \quad  q^{\frac{d-1}{2}}|E|^2 \ll  |E|^{4-4s},\]
which is  equivalent to the condition
$$  q^{\frac{3d-1}{4}} \ll |E|  \quad \mbox{or}\quad q^{\frac{d-1}{4(1-2s)} }\ll |E|.$$
Using  our initial assumption $|E|\ll q^{\frac{d+1}{2}}$ and disregarding the condition $q^{\frac{3d-1}{4}} \ll |E|, $  we take a sufficient condition for the inequality \eqref{Suf*} to hold as follows
\[q^{\frac{d-1}{4(1-2s)} }\ll |E| \ll q^{\frac{d+1}{2}}.\]

For this inequality to hold,  we  also need the condition 
$\frac{d-1}{4(1-2s)} \le \frac{d+1}{2}$, which becomes the condition 
$$\frac{1}{4} \leq s \le \frac{1}{4}+\frac{1}{2(d+1)}.$$
In conclusion,   assumption $(i)$ implies  the $(4,s)$--Salem set property.  Hence, under assumption $(i),$ the proof of Proposition \ref{propro1} is complete.

\textbf{Case 2. Assume that $\frac{|E|^3}{q} \gg q^{\frac{d-1}{2}}|E|^2 (\iff |E| \gg q^{\frac{d+1}{2}})$}.

In this case,   $L\sim \frac{|E|^3}{q}$.   Hence,  the statement \eqref{LR*}  is equivalent to the statement
\begin{equation}\label{ECT*} \frac{|E|^3}{q}  \ll  \frac{|E|^4}{q^d}+ |E|^{4-4s}.\end{equation}

To complete the proof of Proposition \ref{propro1},  it will be enough to show that   each of assumptions $(ii), (iii),$ and $(iv)$ is a sufficient condition for the above inequality to hold.

If $|E|\sim q^{d-1}$ (assumption $(iv))$, then the inequality \eqref{ECT*} obviously holds, because in this case $\frac{|E|^3}{q}\sim \frac{|E|^4}{q^d}$.
Hence,  the proof for the case of assumption $(iv)$ is completed.

To show that each of assumptions $(ii), (iii)$ is a sufficient condition for the inequality \eqref{ECT*} to hold, we need to check the condition $$\frac{|E|^3}{q} \ll |E|^{4-4s},$$
which is the same as 
$$ |E|\ll q^{\frac{1}{4s-1}}.$$
From the original assumption on $|E|$, $|E|\gg q^{\frac{d+1}{2}}$, the sufficient condition takes
$$  q^{\frac{d+1}{2}} \ll |E|\ll q^{\frac{1}{4s-1}}.$$

Since $|E|\ll q^{d-1}$ by \eqref{SizePS},    the sufficient condition  can be taken as
$$  q^{\frac{d+1}{2}} \ll |E|\ll \min\{q^{\frac{1}{4s-1}},  ~~q^{d-1}\}.$$
Observing that 
\[ \min\left\{q^{\frac{1}{4s-1}},  ~~q^{d-1}\right\} = \begin{cases}
    q^{d-1}, \quad                  &\text{if } \frac{1}{4} \le s \le \frac{1}{4}+\frac{1}{4(d-1)},\\
     q^{\frac{1}{4s-1}}, \quad     & \text{if }  \frac{1}{4}+\frac{1}{4(d-1)} \le s \le \frac{1}{4}+ \frac{1}{2(d+1)}, 
    \end{cases}\]
we see that the condition $(ii)$ or $(iii)$ is a sufficient condition for the inequality \eqref{ECT*} as required.
\end{proof}

\end{document}
